\subsection{Editorial verification: what the root-multiplicity question actually asks}

No item in records 530--537 is assigned class C or D. The following proofs and computations are annotations checking the reconstruction and the literal errors; they are not claimed as new research. In particular, an unspecified wish for an illuminating multiplicity formula is not turned into an open demand for a positive formula. The symmetrizable formula in \citet[Theorem 2]{bermanmoody1979multiplicities} was already available in 1979. The elementary argument below covers the finite computational interpretation of 1981-533.4 without importing that paper's symmetrizability hypothesis.

\begin{proposition}[Recursive lowering determines every specified positive degree]\ours
Let \(A\) be a finite generalized Cartan matrix over characteristic zero. Use a Cartan realization with independent simple roots, and define the contragredient algebra by factoring the Chevalley cross-relation algebra by its largest ideal disjoint from the Cartan subalgebra. If \(\alpha\in Q_+\) has height greater than one, then
\[
 \mathfrak g_\alpha\longrightarrow
 \bigoplus_{i:\,\alpha-\alpha_i\in Q_+\setminus\{0\}}
 \mathfrak g_{\alpha-\alpha_i},\qquad
 v\longmapsto([f_i,v])_i
\]
is injective. Consequently a prescribed root-space dimension and independent representing words are computable by finitely many exact linear operations, inductively on height.
\end{proposition}
\begin{proof}
The positive part of the algebra before the maximal quotient is the free Lie algebra on the \(e_i\); its component of degree \(\alpha\) is finite-dimensional. The triangular decomposition and maximal-ideal construction are those of \citet[Chapter II, \S1]{kac1968graded}. The simple positive components have dimension one.

Suppose \(v\in\mathfrak g_\alpha\) is killed by all \(f_i\). Let \(J\) be the span of all iterated brackets
\(
 [e_{i_1},[e_{i_2},\ldots,[e_{i_s},v]\ldots]]
\), including \(v\). It is stable under the \(e_i\) and the Cartan subalgebra. It is stable under each \(f_j\) as well: move \(\operatorname{ad}f_j\) through the iterated positive adjoints, using
\([f_j,e_i]=-\delta_{ij}h_i\), and then move each Cartan adjoint through the remaining positive adjoints. Every resulting term is a scalar multiple of a shorter positive iterated bracket on \(v\); the terminal \([f_j,v]\) vanishes. This induction also covers an arbitrary linear combination of such brackets. Hence \(J\) is an ideal of \(\mathfrak g\).

Every weight of \(J\) has degree \(\alpha+\beta\) for \(\beta\in Q_+\), so \(J\cap\mathfrak h=0\). The maximal-quotient definition prohibits any nonzero such ideal: its preimage would strictly enlarge the largest ideal disjoint from \(\mathfrak h\). Therefore \(v=0\), proving injectivity. Terms with a negative coefficient in \(\alpha-\alpha_i\) are automatically zero; height greater than one excludes a Cartan-valued target. Removing the redundant central directions of Rudakov's enlarged realization does not change these positive components.

For effectiveness, choose a finite basis or spanning list of free Lie words of degree \(\alpha\). By induction suppose quotient coordinate maps are known in every smaller positive degree. The brackets \([f_i,w]\) are calculated from the cross relations and expressed in those already computed quotient coordinates. Stack the resulting matrices. Its kernel is exactly the degree-\(\alpha\) part of the maximal ideal, by the injection just proved. Its rank is the desired dimension, and a set of pivot columns supplies independent word representatives. All coefficients are obtained from the integer entries of \(A\), so exact rational arithmetic suffices for a generalized Cartan matrix. The calculation at every fixed height terminates because all free Lie components in question are finite-dimensional. No complexity bound, canonical choice of pivots or cancellation-free formula is asserted.
\end{proof}

\begin{remark}[Independent check of the lowering argument]\oursComputed
The supplied \texttt{editorial\_checks.py} computes the same quotient dimensions by two different procedures: the lowering-map recursion above, and a degree-by-degree Lie ideal generated by the two explicit Serre relations. Both use an exact free-Lie basis, but the second procedure never uses a lowering operator. For the symmetrizable rank-two pairs \((a,b)=(1,1),(2,1),(3,1),(2,2),(3,2)\), they agree through height eight. The first three height profiles are respectively
\[
 (2,1,0,0,0,0,0,0),\quad
 (2,1,1,0,0,0,0,0),\quad
 (2,1,1,1,1,0,0,0).
\]
These independently recover total dimensions \(3,4,6\). The affine pair \((2,2)\) gives \((2,1,2,1,2,1,2,1)\); the indefinite pair \((3,2)\) gives \((2,1,2,2,3,3,6,7)\). Finite output does not prove infinite-dimensionality; that assertion uses the previously established Cartan classification. The second procedure is used only for symmetrizable rank-two matrices, where the Serre presentation is justified. It is not substituted for the maximal-ideal definition in the general statement above.
\end{remark}

\subsection{Editorial counterexamples for the transformation-group notes}

For 1981-534.1--5 write \(G_{(S)}\) for the pointwise stabilizer and \(G_{\{S\}}\) for the setwise stabilizer. Then the source's closure is
\(
 \operatorname{cl}(S)=\operatorname{Fix}(G_{(S)}).
\)
This distinction is also needed for the restriction homomorphisms in the notebook's fourth axiom. The exact match of those axioms to Seifert's 1984 theorem remains unverified; the following corrections concern elementary statements, independent of that publication-identification gap.

\begin{proposition}[Singleton closure need not mean equal stabilizers]\ours
For every group action,
\[
 \operatorname{cl}\{x\}=\{y:G_x\subseteq G_y\}.
\]
Even in a faithful transitive action the containment can be strict. If the transitive set is finite, it is necessarily equality.
\end{proposition}
\begin{proof}
The displayed equivalence is just the definition that every element fixing \(x\) also fixes \(y\). For strictness let \(G=\mathbb Z[1/2]\rtimes\langle t\rangle\), with \(tbt^{-1}=2b\) on the additive translation subgroup, let \(H=\mathbb Z\), and let \(G\) act on left cosets \(G/H\). At \(x=H\), the stabilizer is \(H\). At \(y=t^{-1}H\), it is \(t^{-1}Ht=\tfrac12\mathbb Z\), strictly larger. Therefore \(y\in\operatorname{cl}\{x\}\) while \(G_y\ne G_x\). The action is transitive by construction and faithful: its kernel lies in \(\bigcap_{n\ge0}2^n\mathbb Z=\{0\}\), hence is trivial. Thus faithfulness does not repair the notebook's statement.

In a finite transitive set, \(G_x\) and \(G_y\) have the same finite index in \(G\). If \(G_x\subseteq G_y\), the index multiplication formula implies \([G_y:G_x]=1\). This is a separate verification of the missing finite-action hypothesis, and does not require \(G\) itself to be finite.
\end{proof}

\begin{example}[The degree-two exception to the double-transitivity remark]\ours
Let \(C_2\) act regularly on its two elements. It is transitive on the two ordered pairs of distinct elements, hence doubly transitive. Every point stabilizer is trivial, so \(\operatorname{cl}\{x\}=X\), strictly larger than \(\{x\}\). Thus the notebook's proposed obstruction requires a restriction. For a doubly transitive set with at least three points, \(G_x\) acts transitively on \(X\setminus\{x\}\), a set of at least two elements, and no point there is fixed by all of \(G_x\). In that range \(\operatorname{cl}\{x\}=\{x\}\) does hold.
\end{example}

\subsection{Independent finite checks of combinatorial formulas}

\begin{proposition}[Interval inversion]\ours
For interval multiplicities \(m_{pq}\) of an equioriented type-A representation,
\[
 r_{ij}=\sum_{p\le i\le j\le q}m_{pq},\qquad
 m_{ij}=r_{ij}-r_{i-1,j}-r_{i,j+1}+r_{i-1,j+1},
\]
where ranks beyond the ends are zero.
\end{proposition}
\begin{proof}
An interval summand contributes one to the rank of the composition from vertex \(i\) to vertex \(j\) exactly when it contains both endpoints, proving the first formula. The difference \(r_{ij}-r_{i-1,j}\) retains exactly intervals starting at \(i\) and ending at or beyond \(j\). Subtracting the corresponding difference with endpoint \(j+1\) retains exactly the intervals ending at \(j\). This proves the second formula. It verifies the indexing in the source completion from \citet[\S2.3]{abeasisdelfrakraft1981geometry}, not the deeper converse degeneration theorem.
\end{proof}

\begin{remark}[Independent enumeration results]\oursComputed
The code constructs explicit direct-sum arrow matrices for all 120 choices of three distinct intervals on four vertices, with independently varied multiplicities one or two. Matrix products and exact ranks agree with the finite-difference reconstruction in every case. For the free Lie formula, direct enumeration of primitive binary necklaces agrees with the M\"obius formula in all 65 bidegrees of total degree one through ten. The independent binary-tree recurrence agrees with the Catalan specialization for \(n=0,\ldots,12\).

For the Gelfand--Tsetlin question, recursive enumeration of every intervening integer row agrees with the Weyl product for all 210 weakly increasing top rows drawn from \(\{-2,-1,0,1,2,3,4\}\). Examples give counts \(1\) for \((0,0,0,0)\), \(4\) for \((0,0,0,1)\), \(64\) for \((0,1,2,3)\), and \(630\) for the additional row \((-2,-1,1,4)\). The general theorem is supplied by the interlacing basis \citep[\S3]{gelfandtsetlin1950unimodular} and the Weyl dimension formula; bounded enumeration is an independent check on indexing, endpoints and normalization.

Finally, multiplying a symmetric-algebra Euler product and separately expanding its formal logarithm through degree fourteen confirms \(C_n=\sum_{d\mid n}d\dim V_d\). This checks that the notebook's rescaled coefficients \(c_k=k\dim V_k\), including the denominator \(k\) in \(p(t)\), are internally consistent. Complete parameter ranges and output are retained beside the code.
\end{remark}

\subsection{Two sign and scaling checks}

\begin{example}[A direct test of the BGG rho shift]\ours
With \([e,f]=h\), \([h,f]=-2f\), let \(M(0)\) be the \(\mathfrak{sl}_2\) Verma module generated by \(v_0\), with \(ev_0=hv_0=0\). The nonzero vector \(fv_0\) has weight \(-2\), and \(e(fv_0)=hv_0=0\). The universal property gives a nonzero map \(M(-2)\to M(0)\); on the polynomial PBW bases it is multiplication by \(f\), so it is injective. Since \(\rho(h)=1\), the weights satisfy \(-2+1=-(0+1)\). They do not satisfy \(-2-1=\pm(0-1)\). This verifies the plus-rho convention against both an explicit module calculation and the shifted definition in \citet[\S3]{bernsteingelfandgelfand1971structure}.
\end{example}

\begin{proposition}[Determinant weights of the supergravity density]\ours
In the determinant-monomial ansatz \(A=|W|^p|\Gamma|^q\), the defining-equation and odd-frame covariance laws used by \citet[\S3, equation (8)]{schwarz1982supergravity} force \(p=4/3\), \(q=-1/3\).
\end{proposition}
\begin{proof}
Replacing the four real defining equations by \(f'=\eta f\) multiplies each determinant by \(\det\eta\), while the product of delta functions has weight \(-1\). Thus \(p+q=1\). Changing the complex odd tangent frame by a \(2\times2\) matrix \(S\) gives \(|W|\) weight \(|\det S|\) and the paired-spinor determinant \(|\Gamma|\) weight \(|\det S|^4\); hence \(p+4q=0\). The two equations have the unique solution stated.

As a separate check, rescale all defining equations by a positive real constant \(c\). Both four-by-four determinants acquire \(c^4\), and the delta product acquires \(c^{-4}\). The final density has total exponent \(16/3-4/3-4=0\), while the preliminary expression \(|\det\Gamma|^{-1/2}\prod\delta(f^k)\) has exponent \(-2-4=-6\). This confirms the source's internal correction. The argument determines exponents within this ansatz; it is not a replacement for the stronger local uniqueness classification attributed to \citet{gaydukromanovschwarz1981democracy}, whose full proof was not accessible here.
\end{proof}
